% File: tikz-hello-world-rlc.tex
% Author: Marcin Szewczyk

% Standalone document for compiling a single TikZ figure
\documentclass{article}

\usepackage[left=0.5cm,right=0.5cm,top=2.5cm,bottom=2.5cm]{geometry}

% Circuitikz with American symbols, European resistor style, straight voltage arrows
\usepackage[american,europeanresistors,RPvoltages,straightvoltages]{circuitikz}

\ctikzset{
	bipoles/thickness=1,   % Line thickness of circuit elements
	label distance=1.4    % Distance of labels from elements
}

% Global font size for all TikZ nodes and labels
\tikzset{font=\large}

% Default arrow tip style (used for current and voltage arrows)
\tikzset{>={Latex[length=5pt,width=5pt]}}

\ctikzset{
	voltage shift=2, % większy odstęp od elementów
	label distance=5  % <-- zwiększa odstęp podpisu elementu od elementu
}

\newlength{\figsep}
\setlength{\figsep}{26pt}

\begin{document}
	
	% Main circuit drawing environment
	% scale controls overall figure size
	% line width sets thickness of wires and elements

\begin{figure}[ht]
	\centering
	\begin{circuitikz}[scale=0.7,line width=1pt]

		\draw
(0,0) to[isource, l=$e(t)$] (0,4)
-- (1.5,4)
to[L, l=$L$, v_<=$u_\mathrm{L}(t)$] (3.5,4)
-- (4.5,4)
to[R, l=$R$, v_<=$u_\mathrm{R}(t)$] (6.3,4)
-- (7.3,4)
to[closing switch, l=$S$, bipoles/length=2cm] (9.3,4)
to[short, i>^=$i(t)$] (11.1,4)
-- (11.1,0) -- (0,0);


	\end{circuitikz}
	\begin{circuitikz}[scale=0.7,line width=1pt]
	
	\draw
(0,0) to[isource, l=$e(t)$] ++(0,4)
-- ++(1.5,0)
to[L, l=$L$, v_<=$u_\mathrm{L}(t)$] ++(2.0,0)
-- ++(1.0,0)
to[R, l=$R$, v_<=$u_\mathrm{R}(t)$] ++(1.8,0)
-- ++(1.0,0)
to[closing switch, l=$S$, bipoles/length=2cm] ++(2.0,0)
to[short, i>^=$i(t)$] ++(1.8,0)
-- ++(0,-4) -- ++(-11.1,0);

	
\end{circuitikz}
	\caption{Circuitikz example: same circuit generated from two different code variants}
\end{figure}

\vspace{\figsep}

\begin{figure}[ht]
	\centering
	\begin{circuitikz}[scale=0.7,line width=1pt]
		
		\newif\ifshowgrid
		\newif\ifshownodes
		
		\showgridfalse
		\shownodestrue
		
		% --- Debug switches ---
% Enable or disable auxiliary drawing elements (grid and nodes)

%\newif\ifshowgrid   % Controls visibility of the background grid
%\newif\ifshownodes  % Controls visibility of debug nodes and labels

%\showgridtrue      % Set to \showgridtrue to show the grid
%\shownodestrue      % Set to \shownodesfalse to hide debug markers and labels

% Debug grid to assist with manual layout alignment
\ifshowgrid
\draw [help lines] (-3,-1) grid (13,6);
\fi

% --- Point definitions ---
\coordinate (n0) at (0,0);
\coordinate (n1) at (0,4);
\coordinate (n2) at (1.5,4);
\coordinate (n3) at (3.5,4);
\coordinate (n4) at (4.5,4);
\coordinate (n5) at (6.3,4);
\coordinate (n6) at (7.3,4);
\coordinate (n7) at (9.3,4);
\coordinate (n8) at (11.1,4);
\coordinate (n9) at (11.1,0);

% Consistent style for debug labels
\tikzset{
	debug label/.style={
		font=\scriptsize,
		fill=white,
		inner sep=0pt
	}
}

% --- Main circuit ---
\draw
(n0) to[isource, l=$e(t)$] (n1)
(n1) to[short, -] (n2)
(n2) to[L=$L$] (n3)
(n3) to[short, -] (n4)
(n4) to[R=$R$] (n5)
(n5) to[short, -] (n6)
(n6) to[/tikz/circuitikz/bipoles/length=2cm, closing switch=$S$] (n7)
(n7) to[short, -, i>^=$i(t)$] (n8)
(n8) to[short, -] (n9)
(n9) to[short, -] (n0)
;

% --- Manually drawn voltage arrows ---
% NOTE: Voltage arrows are drawn manually to control length and alignment

% --- ustaw długość strzałki ---
\def\varrow{1.6}  % całkowita długość

% u_L
\draw[<-]
($(n2)!0.5!(n3)+(-0.5*\varrow,-0.65)$) --
($(n2)!0.5!(n3)+(0.5*\varrow,-0.65)$)
node[midway, below=5pt, fill=white, inner sep=0pt] {$u_\mathrm{L}(t)$};

% u_R
\draw[<-]
($(n4)!0.5!(n5)+(-0.5*\varrow,-0.65)$) --
($(n4)!0.5!(n5)+(0.5*\varrow,-0.65)$)
node[midway, below=5pt, fill=white, inner sep=0pt] {$u_\mathrm{R}(t)$};		
% --- Debug: point names ---
\ifshownodes
\foreach \p in {n1,n2,n3,n4,n5,n6,n7,n8}{
	\fill[red] (\p) circle (2pt);
	\node[above=5pt, debug label] at (\p) {\p};
}
\foreach \p in {n0,n9}{
	\fill[red] (\p) circle (2pt);
	\node[below=5pt, debug label] at (\p) {\p};
}
\fi

		
	\end{circuitikz}
	\caption{Circuitikz example with nodes}
\end{figure}

\vspace{\figsep}
	
\begin{figure}[ht]
	\centering
	\begin{circuitikz}[scale=0.7,line width=1pt]
		
		\newif\ifshowgrid
		\newif\ifshownodes
		
		\showgridtrue
		\shownodestrue
		
		% --- Debug switches ---
% Enable or disable auxiliary drawing elements (grid and nodes)

%\newif\ifshowgrid   % Controls visibility of the background grid
%\newif\ifshownodes  % Controls visibility of debug nodes and labels

%\showgridtrue      % Set to \showgridtrue to show the grid
%\shownodestrue      % Set to \shownodesfalse to hide debug markers and labels

% Debug grid to assist with manual layout alignment
\ifshowgrid
\draw [help lines] (-3,-1) grid (13,6);
\fi

% --- Point definitions ---
\coordinate (n0) at (0,0);
\coordinate (n1) at (0,4);
\coordinate (n2) at (1.5,4);
\coordinate (n3) at (3.5,4);
\coordinate (n4) at (4.5,4);
\coordinate (n5) at (6.3,4);
\coordinate (n6) at (7.3,4);
\coordinate (n7) at (9.3,4);
\coordinate (n8) at (11.1,4);
\coordinate (n9) at (11.1,0);

% Consistent style for debug labels
\tikzset{
	debug label/.style={
		font=\scriptsize,
		fill=white,
		inner sep=0pt
	}
}

% --- Main circuit ---
\draw
(n0) to[isource, l=$e(t)$] (n1)
(n1) to[short, -] (n2)
(n2) to[L=$L$] (n3)
(n3) to[short, -] (n4)
(n4) to[R=$R$] (n5)
(n5) to[short, -] (n6)
(n6) to[/tikz/circuitikz/bipoles/length=2cm, closing switch=$S$] (n7)
(n7) to[short, -, i>^=$i(t)$] (n8)
(n8) to[short, -] (n9)
(n9) to[short, -] (n0)
;

% --- Manually drawn voltage arrows ---
% NOTE: Voltage arrows are drawn manually to control length and alignment

% --- ustaw długość strzałki ---
\def\varrow{1.6}  % całkowita długość

% u_L
\draw[<-]
($(n2)!0.5!(n3)+(-0.5*\varrow,-0.65)$) --
($(n2)!0.5!(n3)+(0.5*\varrow,-0.65)$)
node[midway, below=5pt, fill=white, inner sep=0pt] {$u_\mathrm{L}(t)$};

% u_R
\draw[<-]
($(n4)!0.5!(n5)+(-0.5*\varrow,-0.65)$) --
($(n4)!0.5!(n5)+(0.5*\varrow,-0.65)$)
node[midway, below=5pt, fill=white, inner sep=0pt] {$u_\mathrm{R}(t)$};		
% --- Debug: point names ---
\ifshownodes
\foreach \p in {n1,n2,n3,n4,n5,n6,n7,n8}{
	\fill[red] (\p) circle (2pt);
	\node[above=5pt, debug label] at (\p) {\p};
}
\foreach \p in {n0,n9}{
	\fill[red] (\p) circle (2pt);
	\node[below=5pt, debug label] at (\p) {\p};
}
\fi

		
	\end{circuitikz}
	\caption{Circuitikz example with nodes and grid}
\end{figure}
	
\end{document}
